\chapter{Použití knihovny PlanG}
\label{app:howto}

Tato příloha stručně popisuje, jak sestavit projekt PlanG a jak použít jeho
nejdůležitější části v C++ programu. Úplnější přehled metod je uveden v souboru
\texttt{README.md}; programátorská dokumentace v souboru \texttt{DEVELOPER.md}
popisuje vnitřní návrh knihovny.

\section{Sestavení projektu}

Projekt vyžaduje překladač C++, nástroj
\texttt{make} a knihovnu Boost. Lze jej sestavit na Linuxu a macOS; na Windows
je doporučené použít WSL. Na Linuxu se používá \texttt{gcc} a \texttt{g++},
zatímco na macOS instalační skript použije \texttt{clang} a \texttt{clang++}.
Projekt lze sestavit instalačním skriptem \texttt{install.sh}:

\begin{Verbatim}
./install.sh
\end{Verbatim}

Skript nejprve hledá hlavičkové soubory Boostu v běžných umístěních. Pokud je
nenajde, pokusí se Boost automaticky doinstalovat: na Linuxu pomocí
\texttt{apt-get} a na macOS pomocí Homebrew. Poté sestaví oba backendy,
zkopíruje statické knihovny do adresáře \texttt{lib/} a vytvoří spustitelný
soubor \texttt{plang}. Pokud je knihovna Boost v nestandardním umístění, lze
její cestu předat proměnnou \texttt{BOOST\_DIR}:

\begin{Verbatim}
BOOST_DIR=/path/to/boost ./install.sh
\end{Verbatim}

Výchozí program se spustí příkazem:

\begin{Verbatim}
./plang
\end{Verbatim}

Soubor \texttt{main.cpp} slouží jako hlavní ukázkový program. Po jeho změně je
potřeba projekt znovu sestavit.

\section{Základní program}

Veřejné API knihovny zpřístupňuje hlavičkový soubor \texttt{MyGraphLib.hpp}.
Uživatel zvolí backend pomocí šablony \texttt{Generator<Backend>}. Backend
\texttt{geng::Backend} používá generátor \texttt{geng}, zatímco backend
\texttt{plantri::Backend} používá generátor \texttt{plantri}.

\begin{listing}
\begin{lstlisting}[language=C++]
#include "MyGraphLib.hpp"

using App = Generator<geng::Backend>;

int main()
{
    App app;
    app.setVertices(10);
    app.setConnected();
    app.setTriangleFree();
    return app.run();
}
\end{lstlisting}
\caption{Jednoduché použití backendu \texttt{geng}.}
\label{lst:howto-basic-geng}
\end{listing}

Stejný tvar programu lze použít i pro \texttt{plantri}. Liší se pouze zvolený
backend a metody dostupné pro konkrétní generátor.

\begin{listing}
\begin{lstlisting}[language=C++]
#include "MyGraphLib.hpp"

using App = Generator<plantri::Backend>;

int main()
{
    App app;
    app.setVertices(12);
    app.setFormat(App::OutputFormat::Graph6);
    return app.run();
}
\end{lstlisting}
\caption{Jednoduché použití backendu \texttt{plantri}.}
\label{lst:howto-basic-plantri}
\end{listing}

Oba backendy umožňují vypnout výstup metodou \texttt{setNoOutput} nebo zapsat
výsledek do souboru metodou \texttt{setOutputFile}. Backend \texttt{geng}
podporuje například omezení počtu hran, souvislost, bipartitnost a stupňová
omezení. Backend \texttt{plantri} podporuje volbu třídy rovinných grafů,
minimální stupeň, konektivitu, výstup duálního grafu a různé výstupní formáty.

\section{Callbacky}

Knihovna nabízí tři základní callbacky. Callback \texttt{setFilter} se používá
pro filtrování dokončených grafů. Callback \texttt{setPreprune} se volá během
generování a může předčasně zahodit rozpracovanou větev. Callback
\texttt{setOutproc} slouží k vlastnímu zpracování grafu, který prošel filtrem.

Filtrační callbacky vracejí hodnotu \texttt{KEEP}, pokud má graf nebo větev
zůstat zachována, a hodnotu \texttt{PRUNE}, pokud má být odmítnuta. Následující
program generuje grafy pomocí \texttt{geng} a ponechá pouze grafy, které mají
alespoň jednu hranu.

\begin{listing}
\begin{lstlisting}[language=C++]
#include "MyGraphLib.hpp"

using App = Generator<geng::Backend>;

int main()
{
    App app;
    app.setVertices(6);

    app.setFilter([](const App::GraphView& g) {
        return num_edges(g) > 0 ? KEEP : PRUNE;
    });

    return app.run();
}
\end{lstlisting}
\caption{Filtrování dokončených grafů pomocí \texttt{setFilter}.}
\label{lst:howto-filter}
\end{listing}

Callback \texttt{setOutproc} se hodí tehdy, když program nemá grafy pouze
vypsat v původním formátu, ale chce nad nimi provést vlastní výpočet. V
následujícím příkladu program pouze spočítá, kolik grafů generátor vytvořil.

\begin{listing}
\begin{lstlisting}[language=C++]
#include "MyGraphLib.hpp"
#include <iostream>

using App = Generator<plantri::Backend>;

int main()
{
    App app;
    int count = 0;

    app.setVertices(10);
    app.setNoOutput();

    app.setOutproc([&](Output&, const App::GraphView&) {
        ++count;
    });

    int result = app.run();
    std::cout << count << "\n";
    return result;
}
\end{lstlisting}
\caption{Vlastní zpracování výstupu pomocí \texttt{setOutproc}.}
\label{lst:howto-outproc}
\end{listing}

\section{Práce s grafem}

Callbacky dostávají objekt \texttt{GraphView}. Tento objekt je pohled na graf
uložený uvnitř původního generátoru. Graf tedy není automaticky kopírován do
nové datové struktury. Objekt \texttt{GraphView} je určený pro použití uvnitř
callbacku; po pokračování generování se interní data mohou změnit.

Pro práci s grafem poskytuje PlanG globální funkce, které berou graf jako
parametr. Patří mezi ně například \texttt{vertices(g)}, \texttt{edges(g)},
\texttt{num\_vertices(g)}, \texttt{num\_edges(g)}, \texttt{out\_edges(v, g)},
\texttt{adjacent\_vertices(v, g)}, \texttt{source(e, g)} a
\texttt{target(e, g)}. Tyto funkce lze používat přímo v uživatelském callbacku.

\begin{listing}
\begin{lstlisting}[language=C++]
app.setFilter([](const App::GraphView& g) {
    auto [first, last] = vertices(g);

    for (auto it = first; it != last; ++it)
    {
        int v = *it;
        if (out_degree(v, g) == 0)
            return PRUNE;
    }

    return KEEP;
});
\end{lstlisting}
\caption{Použití globálních funkcí pro práci s objektem \texttt{GraphView}.}
\label{lst:howto-graphview-functions}
\end{listing}

U backendu \texttt{plantri} jsou navíc dostupné funkce pro práci s orientovanými
hranami rovinného vložení, například \texttt{next\_edge},
\texttt{prev\_edge} a \texttt{opposite\_edge}. Tyto funkce jsou užitečné u
algoritmů, které potřebují pracovat přímo s cyklickým pořadím hran kolem
vrcholů.

\section{Použití algoritmů Boost Graph Library}

Objekty \texttt{GraphView} jsou přizpůsobeny vybraným konceptům Boost Graph
Library. Uživatel proto může v callbacku spustit standardní grafový algoritmus,
aniž by graf převáděl do datové struktury z Boostu. Následující příklad používá
algoritmus \texttt{boost::sequential\_vertex\_coloring}.

\begin{listing}
\begin{lstlisting}[language=C++]
#include "MyGraphLib.hpp"
#include <boost/graph/sequential_vertex_coloring.hpp>
#include <vector>

using App = Generator<plantri::Backend>;

int main()
{
    App app;
    app.setVertices(13);

    app.setFilter([](const App::GraphView& g) {
        std::vector<int> colors(num_vertices(g));
        int used = boost::sequential_vertex_coloring(g, colors.data());
        return used >= 4 ? KEEP : PRUNE;
    });

    return app.run();
}
\end{lstlisting}
\caption{Použití algoritmu z Boost Graph Library nad objektem
\texttt{GraphView}.}
\label{lst:howto-boost}
\end{listing}

\section{Obarvené a zakotvené vrcholy v backendu \texttt{geng}}

Backend \texttt{geng} podporuje také generování grafů s barevnými třídami
vrcholů. Přesné velikosti barevných tříd se nastaví metodou
\texttt{setColorClassSizes}. Pokud stačí zadat pouze počet barev, lze použít
\texttt{setColors}. Zakotvené vrcholy se nastavují metodou
\texttt{setRootedVertices}.

\begin{listing}
\begin{lstlisting}[language=C++]
#include "MyGraphLib.hpp"

using App = Generator<geng::Backend>;

int main()
{
    App app;
    app.setVertices(5);
    app.setColorClassSizes({3, 2});
    app.setCanonicalLabeling();

    app.setFilter([](const App::GraphView& g) {
        int cross_edges = 0;

        auto [first, last] = edges(g);
        for (auto it = first; it != last; ++it)
        {
            auto [u, v] = *it;
            if (g.color(u) != g.color(v))
                ++cross_edges;
        }

        return cross_edges >= 2 ? KEEP : PRUNE;
    });

    return app.run();
}
\end{lstlisting}
\caption{Generování grafů s dvěma barevnými třídami.}
\label{lst:howto-colored}
\end{listing}

Barvy lze v callbacku číst přímo z objektu \texttt{GraphView}. K dispozici jsou
metody \texttt{has\_coloring()}, \texttt{color\_count()},
\texttt{color(v)}, \texttt{vertex\_colors()} a
\texttt{vertices\_of\_color(c)}. Zakotvené vrcholy jsou reprezentovány jako
samostatné jednoprvkové barevné třídy.

\section{Dostupné ukázky a benchmarky}

Adresář \texttt{examples/} obsahuje další ukázkové programy demonstrující použití knihovny.
Adresář \texttt{benchmarks/} obsahuje programy použité pro měření v této práci.



Zdrojové kódy projektu jsou dostupné v repozitáři: \url{https://github.com/dochynes/PlanG/tree/main}